The primary comparison is a fixed-replay experiment designed to isolate resource
allocation. It replays $76$ held-out routes from the temporal test partition,
uses the same integrated simulation environment and the same case arrivals for
every method. The BPMN model is loaded for candidate generation and diagnostics,
but the fixed replay continues along the historical route even when a diagnostic
BPMN firing fails. Therefore the reported completion value is a fixed-route
completion rate, not a BPMN final-marking rate.

The corrected canonical run uses the leakage-free RF-optimized BPMN-replay evaluation
artifact
\texttt{joao/models/branching/composite\_branching\_evaluation\_train70\_rfopt\_v1.pkl}
(SHA-256:
\texttt{490da841440cd019a0819892ab659e5b61f5ee55182f6db12a536a9c4d23ff82}),
the train-split ParkSong processing-time estimates, seeds $1$--$5$, and the
five methods RoundRobin, ShortestQueue, ParkSong-Composite,
Kunkler-Rinderle-Ma and Batch. The arrival window is 2016-09-16 15:07 UTC to
2016-09-17 15:07 UTC, followed by a drain horizon until 2016-11-16 15:07 UTC.
Cycle time is computed only for completed routes. Waiting time measures resource
queue entry to resource assignment. Throughput and occupation are normalized by
the arrival-plus-drain horizon, so they are retained mainly in the appendix.
Because cycle time is conditional on completed routes, censored route instances
are reported separately and the cycle-time ranking should be read together with
the completion rate. Waiting-time summaries likewise describe tasks that reached
a resource-assignment timestamp; they are not an imputed waiting time for tasks
that remained unresolved at the horizon.

A second run is reported only as an integration smoke. It uses a three-hour
generative arrival window, an eight-hour drain window and seeds $1$--$3$. This
run uses the transition-aware BPMN path with an empty composite branching
hierarchy, so it checks transition-aware BPMN and allocation integration rather
than the offline Random-Forest composite classifier. Consequently, the
generative results are not used to rank allocation methods and do not claim
end-to-end Random-Forest composite branching performance.

The corrected final tables are read from
\texttt{joao/results/final\_canonical\_rfopt\_candidate\_20260717/fixed\_replay/}.
Older canonical packages and the full-log
\texttt{final\_composite\_branching.pkl} artifact are retained only as
historical material and are not used for the corrected final tables.
